\chapter{Sackur--Tetrode Equation}

\section*{Introduction}

In classical statistical mechanics, the entropy of an ideal gas contains an arbitrary additive constant---the entropy is only defined up to a reference point. The Sackur--Tetrode equation eliminates this ambiguity by using quantum mechanics to count the number of accessible microstates properly. It maps the macroscopic thermodynamic state $(N, V, T)$ onto an absolute measure of entropy, confirming that quantum mechanics is necessary even for understanding classical gases.
\section*{Fundamental Equation Used}

\begin{equation}
S = Nk\ln\left[\frac{V}{N}\left(\frac{4\pi m E}{3Nh^2}\right)^{3/2}\right] + \frac{5Nk}{2}
\end{equation}
\section*{Assumptions}

\textbf{Assumptions:} Monatomic ideal gas; particles are indistinguishable; quantum statistics (Boltzmann regime, $n_Q \gg n$, where $n_Q$ is the quantum concentration).


\textbf{Limitations:} Does not apply to diatomic or polyatomic gases (rotational and vibrational degrees of freedom omitted); breaks down at very low temperatures where quantum degeneracy becomes important.


\section*{Mathematical Derivation}

\textbf{Step 1: Write the single-particle partition function for a monatomic ideal gas.}

For a single particle of mass $m$ in a volume $V$ at temperature $T$, the translational partition function is
\begin{equation}
Z_1 = \frac{1}{h^3}\int e^{-\beta p^2/2m}\,d^3\mathbf{r}\,d^3\mathbf{p} = \frac{V}{h^3}\left(\frac{2\pi m}{\beta}\right)^{3/2} = \frac{V}{\lambda^3},
\end{equation}
where $\lambda = h/\sqrt{2\pi m k_B T}$ is the thermal de Broglie wavelength and $\beta = 1/(k_BT)$.

\textbf{Step 2: Construct the $N$-particle partition function for indistinguishable particles.}

For $N$ identical, indistinguishable particles with no inter-particle interactions:
\begin{equation}
Z_N = \frac{Z_1^N}{N!} = \frac{1}{N!}\left(\frac{V}{\lambda^3}\right)^N.
\end{equation}
The $1/N!$ factor corrects for the overcounting of microstates due to particle indistinguishability (Gibbs factor).

\textbf{Step 3: Compute the Helmholtz free energy.}

\begin{equation}
F = -k_BT\ln Z_N = -k_BT\left[N\ln\frac{V}{\lambda^3} - \ln N!\right].
\end{equation}
Using Stirling's approximation $\ln N! \approx N\ln N - N$:
\begin{equation}
F = -Nk_BT\left[\ln\frac{V}{N\lambda^3} + 1\right].
\end{equation}

\textbf{Step 4: Derive the entropy.}

The entropy is $S = -(\partial F/\partial T)_{N,V}$:
\begin{equation}
S = Nk_B\ln\frac{V}{N\lambda^3} + Nk_BT\frac{\partial}{\partial T}(3\ln\lambda) + Nk_B.
\end{equation}
Since $\lambda \propto T^{-1/2}$, we have $3\ln\lambda = \frac{3}{2}\ln T + \text{const}$, and
\begin{equation}
S = Nk_B\ln\frac{V}{N\lambda^3} + \frac{3Nk_B}{2} + Nk_B.
\end{equation}

\textbf{Step 5: Substitute $\lambda$ and express in terms of energy.}

Using $\lambda = h/\sqrt{2\pi m k_BT}$ and $E = \frac{3}{2}Nk_BT$:
\begin{equation}
\ln\frac{V}{N\lambda^3} = \ln\left[\frac{V}{N}\left(\frac{2\pi mk_BT}{h^2}\right)^{3/2}\right] = \ln\left[\frac{V}{N}\left(\frac{4\pi m E}{3Nh^2}\right)^{3/2}\right].
\end{equation}
Thus:
\begin{equation}
\boxed{S = Nk_B\ln\left[\frac{V}{N}\left(\frac{4\pi m E}{3Nh^2}\right)^{3/2}\right] + \frac{5Nk_B}{2}.}
\end{equation}
The $5Nk_B/2$ arises from $\frac{3}{2}Nk_B$ (kinetic energy entropy) plus $Nk_B$ (spatial mixing entropy from the $1/N!$ correction).

\section*{Characteristics of the Equation}

\begin{itemize}
    \item The equation is extensive: $S$ scales linearly with $N$ (the $V/N$ and $E/N$ ratios ensure this).
    \item The $5Nk/2$ term arises from the combination of translational, spatial, and Gibbs corrections.
    \item The $h^3$ in the denominator is the quantum of phase-space volume, making the argument of the logarithm dimensionless.
    \item Solved the Gibbs paradox for indistinguishable particles; published by Sackur (1911) and Tetrode (1912).
\end{itemize}
\section*{Solution Methods}

\begin{itemize}
    \item \textbf{From the partition function:} Compute $Z_1 = V/\lambda^3$ for a single particle, then $Z_N = Z_1^N/N!$, and derive $S = k\ln Z + kT(\partial\ln Z/\partial T)_{N,V}$.
    \item \textbf{From the energy:} Use $E = \frac{3}{2}NkT$ to eliminate $T$ in favor of $E$.
    \item \textbf{Direct count:} Count the number of quantum states in the volume $h^3$ of phase space accessible to $N$ particles.
\end{itemize}
\section*{Verifications}

\begin{itemize}
    \item The predicted entropies of He, Ne, and Ar agree with experimental values to within a few percent.
    \item The Gibbs paradox is resolved: swapping identical particles produces no entropy change, unlike swapping distinguishable particles.
    \item Dimensional analysis confirms that the argument of the logarithm is dimensionless.
\end{itemize}
\section*{Applications}

\begin{itemize}
    \item Calculating the absolute entropy of noble gases (He, Ne, Ar) for comparison with calorimetric measurements.
    \item Verifying the third law of thermodynamics for ideal systems.
    \item Estimating evaporation entropy and boiling points of simple liquids.
    \item Foundation for more complex statistical mechanical calculations.
\end{itemize}
\section*{What Richard Feynman Would Have Asked}

\subsection*{Plain-English Translation}
\textit{``What is the Sackur--Tetrode equation telling us about nature that we didn't already know from classical thermodynamics?''}

It tells us that entropy has an absolute value, not just relative differences. Classical thermodynamics can only measure changes in entropy, not the entropy itself. The Sackur--Tetrode equation gives the actual number of microstates (through $S = k\ln\Omega$) for a monatomic gas, and it agrees perfectly with experiment. This was one of the first triumphs of statistical mechanics, proving that the microscopic count of states matches macroscopic measurements.

\subsection*{Localized Mechanics}
\textit{``What is happening at the level of individual gas molecules that gives rise to the entropy in the Sackur--Tetrode equation?''}

Each molecule occupies a volume of phase space (position and momentum) of order $h^3$. The entropy counts how many of these phase-space cells the gas can access. More volume means more positions available; more energy means more momenta accessible; and the $V/N$ ratio accounts for the fact that each particle only ``sees'' its share of the total volume (due to indistinguishability). The entropy is the logarithm of this count, and the Sackur--Tetrode equation is the exact result for an ideal gas.

\subsection*{Testing Extreme Limits}
\textit{``What happens to the Sackur--Tetrode equation in the limit of extremely high temperature or extremely large volume?''}

In both limits, the entropy increases without bound, as expected: $S \propto \ln V$ at fixed $E$, and $S \propto \ln E$ at fixed $V$. The argument of the logarithm grows, reflecting the increasing number of accessible microstates. In the opposite limit of very low temperature or very small volume, the equation eventually breaks down because the quantum degeneracy condition $n \ll n_Q$ is violated, and the gas enters a regime where quantum statistics (Bose-Einstein or Fermi-Dirac) must be used.

\subsection*{Dimensionality and Geometry}
\textit{``How would the Sackur--Tetrode equation change if we lived in two spatial dimensions instead of three?''}

In two dimensions, the $3/2$ power becomes $1$ (since the density of states in momentum space scales as $p\,dp$ rather than $p^2\,dp$), and the coefficient changes from $5Nk/2$ to $2Nk$. The general structure remains the same---the entropy is still $Nk$ times a logarithm of a dimensionless quantity---but the specific numbers change because the number of degrees of freedom changes. The $h^2$ replaces $h^3$ as the phase-space volume per particle in 2D.

\subsection*{Cross-Disciplinary Connection}
\textit{``Where in information theory or other fields does a similar logarithmic counting formula appear?''}

The Sackur--Tetrode equation is structurally identical to the Shannon entropy formula $H = -\sum p_i \ln p_i$ in information theory, and to the Boltzmann entropy $S = k\ln\Omega$ itself. In information theory, the entropy counts the number of bits needed to specify a message; in statistical mechanics, it counts the number of microstates needed to specify a macrostate. The logarithmic structure is universal---it appears whenever you count the number of ways something can be arranged, and it ensures that independent systems have additive entropies.

% ------------------------------------------------------------
\section*{FAQ}

\textbf{Q: What is the Gibbs paradox?}
A: If you treat identical particles as distinguishable, the entropy of mixing two identical gases at the same temperature and pressure is positive, which is unphysical. The $1/N!$ factor in the partition function (which leads to the $V/N$ in the Sackur--Tetrode equation) resolves this by accounting for particle indistinguishability.

\textbf{Q: Why is $h$ (Planck's constant) appearing in a classical thermodynamics result?}
A: The entropy depends on how many microstates correspond to a macrostate. The number of microstates depends on the fundamental unit of phase space, which is $h^3$ per particle in three dimensions. This is a quantum mechanical result---classical mechanics cannot determine the absolute number of microstates.

\textbf{Q: What happens at very low temperatures?}
A: At very low temperatures, the gas enters the quantum degenerate regime (Bose-Einstein or Fermi-Dirac statistics), and the Sackur--Tetrode equation is no longer valid. The gas becomes a Bose or Fermi condensate, requiring the full quantum statistical treatment.
\section*{Important Bibliography}

\begin{enumerate}
\item Pathria, R.K. and Beale, P.D., \textit{Statistical Mechanics}, 4th ed. (Academic Press, 2021), Chapter 1.
\item Kittel, C. and Kroemer, H., \textit{Thermal Physics}, 2nd ed. (W.H.\ Freeman, 1980), Chapter 8.
\item Sackur, O., ``Die Grundgleichungen der statistischen Mechanik,'' \textit{Ann.\ Phys.\ } \textbf{345}, 67--86 (1913).
\item Tetrode, H., ``Die chemische Konstante der Gase und das elektrische Elementarquantum,'' \textit{Ann.\ Phys.\ } \textbf{348}, 497--522 (1912).
\end{enumerate}


